% No numbered results (statements-of-record rows resolved here: none).
\section{Discussion}
\label{sec:discussion}

\paragraph{A worked instance: self-adjointness.}
The definitions have a transparent instance in matrix analysis.  Let
\(X\) be the space \(M_d(\mathbb C)\) of complex \(d\times d\) matrices,
\(\mu\) a finite positive Borel measure on \(X\) with
\(\int_X\|T\|_{\mathrm{HS}}^2\,d\mu(T)<\infty\), and \(J(T)=T^*\).
Regard \(Y=M_d(\mathbb C)\) as a real Hilbert space with the inner
product \(\langle A,B\rangle=\operatorname{Re}\operatorname{tr}(A^*B)\),
let \(\Jiso(A)=A^*\), which is a real-linear isometric involution, and
let \(\psi\) be the identity.  Then
\[
  \psim(T)=\tfrac12\,(T-T^*)
\]
is the skew-Hermitian part of \(T\), the fixed locus is the set of
Hermitian matrices, and \(\psim\) is qualitatively separating.  Since
the Hermitian part \(\tfrac12(T+T^*)\) is the Hermitian matrix nearest
to \(T\) in the Hilbert--Schmidt norm, \(\dist(T,\Fix(J))=\|\psim(T)\|\)
and \(\psim\) is quantitatively separating with \(m(\varepsilon)=\varepsilon\).
\Cref{thm:duality_confinement:master-theorem} then says: \(\mu\) is
carried by the Hermitian matrices if and only if the Gram operator of
the skew-Hermitian parts is dominated by positive operators of
vanishing trace.  \Cref{thm:duality_confinement:quantitative-confinement}
bounds the \(\mu\)-mass of matrices at distance at least \(\varepsilon\)
from the Hermitian ones by \(\varepsilon^{-2}\trace\AX\).  Since
\(\trace\AX=\int_X\|\psim\|^2\,d\mu\) can be computed directly here,
the instance illustrates the definitions but gains nothing from the
domination formulation; a substantive use would need a situation in
which \(\AX\) is known only through a bridge to independently computed
operators.

\paragraph{Curves over finite fields.}
For a smooth projective curve over a finite field, the analogue of the
Riemann hypothesis is a theorem of Weil \citep{Weil1948a,Weil1948b}.
Stepanov introduced an elementary polynomial method and used it to
bound the number of points on hyperelliptic curves over prime fields
\citep{Stepanov1969}; Bombieri, building on that method, gave an
elementary proof of Weil's theorem for all curves \citep{Bombieri1974}.
The zeta function of a curve of genus \(g\) over \(\mathbb F_q\) has
\(2g\) inverse roots \(\alpha\), permuted by the involution
\(\alpha\mapsto q/\overline{\alpha}\) coming from the functional
equation; its fixed locus is the circle \(|\alpha|=\sqrt q\).  The
corresponding ledger is therefore a finite weighted one.  Since the
conclusion is known, trace-vanishing records exist trivially (the zero
sequence), so the meaningful question is a different one: whether the
Stepanov--Bombieri counting argument can be recast as a bridge
\(\AX\preceq\Kminus+E\) in which \(\Kminus\) and \(E\) are computed
from point counts rather than from the roots.  We have not carried
this out.

\paragraph{Relation to classical operator theory.}
Stripped of the framework, \Cref{thm:duality_confinement:master-theorem}
is the familiar fact that a positive trace-class operator dominated by
operators of arbitrarily small trace is zero, applied to the Gram
operator of an odd observable.  What the formulation adds is a fixed
interface: an application must specify the involution, the equivariant
readout and the domination records, and the conclusion is then read
off without further analysis.

\paragraph{Further directions.}
The natural direction is to find instances in which domination
records can be produced by estimates, for instance through the
factorization criterion of \Cref{sec:domination}.
